\documentclass[11pt]{article}
\usepackage[margin=1in]{geometry}
\usepackage{amsmath,amssymb,amsthm,mathtools}
\usepackage{graphicx}
\usepackage{microtype}
\usepackage[colorlinks=true,linkcolor=blue,urlcolor=blue,citecolor=blue]{hyperref}

\newtheorem{theorem}{Theorem}
\newcommand{\R}{\mathbb R}
\newcommand{\C}{\mathbb C}
\newcommand{\F}{\mathcal F}
\newcommand{\Dom}{\operatorname{Dom}}

\title{The proposed upper pseudomode curve is not optimal\\
for the linear imaginary biharmonic potential}
\author{Candidate counterexample to the optimality question in arXiv:2201.03305}
\date{11 August 2026}

\begin{document}
\maketitle

\noindent\fbox{\parbox{0.94\linewidth}{\textbf{Status: candidate full
negative answer to the universal/full-line optimality claim, likely valid.}
For
\[
                  H=\frac{d^4}{dx^4}+ix
                  \quad\text{on }L^2(\R),
\]
there are pseudomodes at every $\lambda=\alpha+i\beta$ with
$\alpha\to+\infty$, uniformly in $\beta\in\R$.  Their residual is at most
\[
 C_\varepsilon\alpha^{1/16}
 \exp\!\left[-\left(\frac85-\varepsilon\right)\alpha^{5/4}\right].
\]
In particular, $\alpha=\beta^3$ gives pseudomodes far beyond the source's
upper curve $\alpha\lesssim\beta^{8/3-o(1)}$ for $V(x)=ix$.  Thus that curve
is not an optimal boundary for the pseudospectrum of the original full-line
operator.}}

\section{The source question}

Nguyen Duc studies the maximal full-line operator
\[
             H_V=\frac{d^4}{dx^4}+V(x)
\]
and constructs pseudomodes near a positive-end turning point
$x_\beta$, defined by $\operatorname{Im}V(x_\beta)=\beta$.  Theorem 2.2
assumes, among other conditions,
\begin{equation}\label{eq:sourcecurve}
 [\beta\tau(x_\beta)]^{4/5}\lesssim |\alpha|
 \lesssim
 [\beta\tau(x_\beta)^{-1}]^{(4+\varepsilon_1)/(3+\varepsilon_1)}.
\end{equation}
Immediately afterward, the paper states that when
$\operatorname{Re}V=0$ it expects both bound curves in
\eqref{eq:sourcecurve} to be optimal and calls their optimality an open
problem \cite[p.~7, Theorem~2.2 and the following paragraph]{NguyenDuc}.
Here is the complete source context.

\begin{center}
\includegraphics[width=\linewidth]{figures/open_problem_crop.png}
\end{center}

We show that the proposed upper curve is not a necessary boundary, even for
the admissible purely imaginary polynomial $V(x)=ix$, provided the question
is read for the original full-line operator.

\section{Proof intuition}

The linear imaginary potential has an exact symmetry that the local WKB
analysis does not use.  Fourier transform sends multiplication by $ix$ to
$-\partial_\xi$, so $H-(\alpha+i\beta)$ becomes
\[
             -\partial_\xi+\xi^4-\alpha-i\beta.
\]
The phase $e^{-i\beta\xi}$ removes $\beta$ exactly.  For
$\alpha=a^4$, the remaining first-order equation has the exact homogeneous
solution
\[
                  e^{\Phi_a(\xi)},\qquad
                  \Phi_a(\xi)=\frac{\xi^5}{5}-a^4\xi.
\]
Its exponent has a local maximum $\frac45a^5$ at $\xi=-a$ and a local
minimum $-\frac45a^5$ at $\xi=a$.  We retain the maximum and cut the solution
off just after the minimum.  The norm is controlled by the maximum, while the
cutoff error is controlled by the minimum, producing the barrier
$\frac85a^5=\frac85\alpha^{5/4}$.  This construction is uniform in the
imaginary spectral coordinate.

\section{Uniform full-line pseudomodes}

Let $H$ be the maximal operator associated with $u\mapsto u^{(4)}+ixu$ on
$L^2(\R)$.  Every Schwartz function belongs to $\Dom(H)$.

\begin{theorem}\label{thm:main}
For every $\varepsilon>0$ there are constants $C_\varepsilon>0$ and
$A_\varepsilon>0$ such that, for every $\alpha\ge A_\varepsilon$ and every
$\beta\in\R$, there is a nonzero $u_{\alpha,\beta}\in\mathcal S(\R)$ with
\begin{equation}\label{eq:residual}
 \frac{\|(H-(\alpha+i\beta))u_{\alpha,\beta}\|_2}
      {\|u_{\alpha,\beta}\|_2}
 \le C_\varepsilon\alpha^{1/16}
 \exp\!\left[-\left(\frac85-\varepsilon\right)\alpha^{5/4}\right].
\end{equation}
Consequently, if $\alpha_n\to\infty$, then
$\alpha_n+i\beta_n$ has pseudomodes for every real sequence $(\beta_n)$.
\end{theorem}

\begin{proof}
Use the unitary Fourier transform with
$\widehat{xf}=i\partial_\xi\widehat f$.  Then
\begin{equation}\label{eq:fourierop}
 \F(H-(\alpha+i\beta))\F^{-1}
 =-\partial_\xi+\xi^4-\alpha-i\beta.
\end{equation}
Put $a=\alpha^{1/4}$ and
\begin{equation}\label{eq:phase}
 \Phi_a(\xi)=\frac{\xi^5}{5}-a^4\xi
             =a^5F(\xi/a),
 \qquad
 F(t)=\frac{t^5}{5}-t.
\end{equation}
Thus $(e^{\Phi_a})'=(\xi^4-\alpha)e^{\Phi_a}$.

Fix $\varepsilon>0$.  Since $F(1)=-4/5$, choose $\delta>0$ so small that
\begin{equation}\label{eq:delta}
                    F(1+\delta)\le -\frac45+\varepsilon.
\end{equation}
Choose $\chi\in C_c^\infty((-3,1+\delta))$ with
$\chi=1$ on $[-2,1]$, and set
\begin{equation}\label{eq:v}
 v_{a,\beta}(\xi)
 =e^{-i\beta\xi}\chi(\xi/a)e^{\Phi_a(\xi)}.
\end{equation}
This is a nonzero compactly supported smooth function, so
$u_{\alpha,\beta}=\F^{-1}v_{a,\beta}$ is Schwartz.  Substitution into
\eqref{eq:fourierop} gives the exact cancellation
\begin{equation}\label{eq:error}
 \bigl(-\partial_\xi+\xi^4-\alpha-i\beta\bigr)v_{a,\beta}
 =-a^{-1}e^{-i\beta\xi}\chi'(\xi/a)e^{\Phi_a(\xi)}.
\end{equation}

We now compare the two sides.  The phase has
\[
 F(-1)=\frac45,\qquad F'(-1)=0,\qquad F''(-1)=-4.
\]
On the interval $|\xi+a|\le a^{-3/2}$, Taylor's theorem therefore gives
$\Phi_a(\xi)\ge\frac45a^5-C$ for all large $a$.  This interval lies where
$\chi(\xi/a)=1$, hence
\begin{equation}\label{eq:denom}
             \|v_{a,\beta}\|_2
             \ge c a^{-3/4}e^{\frac45a^5}.
\end{equation}
The constants do not depend on $\beta$ because the extra phase has modulus
one.

The support of $\chi'$ is contained in
$(-3,-2)\cup(1,1+\delta)$.  The function $F$ is increasing on each of these
intervals, and
\[
 F(-2)=-\frac{22}{5},\qquad
 F(1+\delta)\le-\frac45+\varepsilon.
\]
It follows from \eqref{eq:error}, after accounting for a support of length
$O(a)$, that
\begin{equation}\label{eq:numer}
 \|(-\partial_\xi+\xi^4-\alpha-i\beta)v_{a,\beta}\|_2
 \le C_\varepsilon a^{-1/2}
       e^{(-\frac45+\varepsilon)a^5}.
\end{equation}
Dividing \eqref{eq:numer} by \eqref{eq:denom}, using Plancherel, and recalling
$a=\alpha^{1/4}$ yields
\[
 \frac{\|(H-(\alpha+i\beta))u_{\alpha,\beta}\|_2}
      {\|u_{\alpha,\beta}\|_2}
 \le C_\varepsilon a^{1/4}
      e^{-(\frac85-\varepsilon)a^5}
 =C_\varepsilon\alpha^{1/16}
      e^{-(\frac85-\varepsilon)\alpha^{5/4}},
\]
which is \eqref{eq:residual}.  The final assertion is immediate.
\end{proof}

\section{Failure of the proposed upper curve}

For $V(x)=ix$, the source's hypotheses hold with
\[
 x_\beta=\beta,\qquad
 \tau(x)=(1+x^2)^{-1/2},\qquad
 \tau(x_\beta)\asymp\beta^{-1}.
\]
Thus the upper curve in \eqref{eq:sourcecurve} has order
\begin{equation}\label{eq:upper}
 [\beta\tau(x_\beta)^{-1}]^{(4+\varepsilon_1)/(3+\varepsilon_1)}
 \asymp
 \beta^{2(4+\varepsilon_1)/(3+\varepsilon_1)},
\end{equation}
whose exponent is strictly less than $8/3$.  Take
$\alpha=\beta^3$.  Then \eqref{eq:upper} is violated by a power of $\beta$,
whereas Theorem~\ref{thm:main} gives
\[
 \frac{\|(H-(\beta^3+i\beta))u_{\beta^3,\beta}\|_2}
      {\|u_{\beta^3,\beta}\|_2}
 \le C_\varepsilon\beta^{3/16}
       \exp\!\left[-\left(\frac85-\varepsilon\right)\beta^{15/4}\right]
 \longrightarrow0.
\]
Therefore the proposed upper curve is not an optimal boundary for the
full-line pseudospectrum, even in the purely imaginary linear case.

\section{Checks, scope, and novelty}

The argument is exact up to the elementary cutoff estimates: the Fourier
identity, removal of $\beta$, and residual formula \eqref{eq:error} contain no
asymptotic expansion.  The constructed Fourier functions are in
$C_c^\infty$, so their inverse transforms lie in the maximal operator domain.
The exponent $8/5$ is the height difference between the local maximum and
minimum of the phase and can be approached arbitrarily closely.

This result answers the universal/full-line optimality belief negatively.  It
does \emph{not} decide whether the same upper curve is optimal for a genuine
half-line realization with a boundary condition at $0$, or within the
narrower class of pseudomodes required to be supported entirely in
$\R^+$.  Indeed, the source's positive-end WKB construction imposes precisely
such localization, while the full-line symmetry is what removes $\beta$ here.
The lower curve in \eqref{eq:sourcecurve} is also not addressed.

\paragraph{Novelty check.}
Bounded searches through 11 August 2026 covered the exact source title and
optimality wording, ``biharmonic'' or ``fourth-order complex Airy'' with
``pseudomode'' or ``resolvent'', and the model $d^4/dx^4+ix$.  No later paper
recording this Fourier-space counterexample was found.  Searches did find
the well-studied second-order complex Airy operator and later work on complex
magnetic Laplacians, but neither answers this fourth-order curve-optimality
question.  Novelty confidence is moderate because the search was bounded.

\paragraph{Human-review recommendation.}
Check the Fourier sign in \eqref{eq:fourierop}, the exact phase cancellation
in \eqref{eq:error}, the $a^{-3/4}$ norm scale near $\xi=-a$, and the
interpretation of the source's optimality sentence as a claim about the
full-line pseudospectrum rather than only positive-half-line-supported WKB
pseudomodes.

\begin{thebibliography}{9}
\bibitem{NguyenDuc}
T.~Nguyen Duc,
\emph{Pseudomodes for biharmonic operators with complex potentials},
Asymptotic Analysis 132 (2023), no.~3--4, 451--493;
arXiv:2201.03305,
\url{https://arxiv.org/abs/2201.03305}.
\end{thebibliography}

\end{document}
